% =====================================================================
%  00_preamble.tex – ustawienia wyglądu zgodne z Zasadami dyplomowania
%  WEAIiIB AGH, rozdz. XI "Założenia edytorskie"
% =====================================================================

% --- Kodowanie i język polski ---------------------------------------
\usepackage[utf8]{inputenc}
\usepackage[T1]{fontenc}
\usepackage[polish]{babel}

% --- Czcionka: Times New Roman 12 pt (wymóg) ------------------------
% newtxtext/newtxmath = Times dla tekstu i wzorów matematyczny-ch
\usepackage{newtxtext}
\usepackage{newtxmath}
\usepackage{amsmath}
% --- Marginesy (wymóg) -----------------------------------------------
% strony nieparzyste: lewy 3 cm, pozostałe 2 cm
% strony parzyste:    prawy 3 cm, pozostałe 2 cm
% przy opcji twoside: inner = margines od strony grzbietu
\usepackage[a4paper,
            inner=3cm, outer=2cm,
            top=2cm,   bottom=2cm,
            footskip=1cm]{geometry}

% --- Odstęp między wierszami: 1,25 wiersza (wymóg) -------------------
% W Wordzie "1,25 wiersza" dla Times 12 pt to ok. 17,3 pt między liniami.
% LaTeX dla 12 pt ma domyślnie 14,5 pt, więc współczynnik 1.2 daje ten sam efekt.
\usepackage{setspace}
\setstretch{1.2}

% --- Akapity ---------------------------------------------------------
\usepackage{indentfirst}          % wcięcie także w 1. akapicie (polski zwyczaj)
\setlength{\parindent}{1.25cm}
\usepackage{microtype}            % ładniejsze dzielenie i justowanie

% --- Nagłówki rozdziałów: "1. Wstęp" zamiast "Rozdział 1" -----------
\usepackage{titlesec}
\titleformat{\chapter}[hang]{\normalfont\LARGE\bfseries}{\thechapter.}{0.5em}{}
\titlespacing*{\chapter}{0pt}{0pt}{1.5em}
\titleformat{\section}{\normalfont\Large\bfseries}{\thesection.}{0.5em}{}
\titleformat{\subsection}{\normalfont\large\bfseries}{\thesubsection.}{0.5em}{}

% --- Numery stron: na dole, po zewnętrznej stronie --------------------
\usepackage{fancyhdr}
\pagestyle{fancy}
\fancyhf{}
\fancyfoot[LE,RO]{\thepage}
\renewcommand{\headrulewidth}{0pt}
\fancypagestyle{plain}{\fancyhf{}\fancyfoot[LE,RO]{\thepage}%
  \renewcommand{\headrulewidth}{0pt}}  % strony z początkiem rozdziału
\usepackage{emptypage}           % puste strony wstawiane przez \cleardoublepage – bez numeru

% --- Grafika, tabele, podpisy ----------------------------------------
\usepackage{graphicx}
\graphicspath{{img/}}             % obrazki szukane w katalogu img/
\usepackage{booktabs}             % ładne tabele (\toprule, \midrule ...)
\usepackage{float}                % opcja [H] – obrazek dokładnie tutaj
\usepackage[font=small,labelfont=bf,labelsep=period]{caption}
\usepackage{siunitx}              % jednostki: \SI{1090}{\mega\hertz}
\sisetup{output-decimal-marker={,}}
\usepackage{enumitem}             % kontrola wyglądu list
% Listy (punktowane i numerowane): znak listy na wysokości wcięcia
% akapitu (1,25 cm), tekst zaraz za nim - jednakowo w całej pracy
\setlist[itemize,enumerate]{labelindent=\parindent, leftmargin=*}
\setlist{itemsep=0pt}

% --- Kod źródłowy (Python, bash, YAML) -------------------------------
\usepackage{xcolor}
\usepackage{listings}
\lstset{
  basicstyle=\ttfamily\small\linespread{1.0}\selectfont,
  numbers=left, numberstyle=\tiny, numbersep=8pt,
  xleftmargin=2em, framexleftmargin=1.5em,
  frame=single, breaklines=true, columns=fullflexible,
  keywordstyle=\color{blue!60!black}\bfseries,
  commentstyle=\color{green!40!black}\itshape,
  stringstyle=\color{red!60!black},
  showstringspaces=false, tabsize=4,
  captionpos=b,
  % polskie znaki w kodzie (pdfLaTeX):
  literate={ą}{{\k{a}}}1 {ć}{{\'c}}1 {ę}{{\k{e}}}1 {ł}{{\l{}}}1
           {ń}{{\'n}}1 {ó}{{\'o}}1 {ś}{{\'s}}1 {ź}{{\'z}}1 {ż}{{\.z}}1
           {Ą}{{\k{A}}}1 {Ć}{{\'C}}1 {Ę}{{\k{E}}}1 {Ł}{{\L{}}}1
           {Ń}{{\'N}}1 {Ó}{{\'O}}1 {Ś}{{\'S}}1 {Ź}{{\'Z}}1 {Ż}{{\.Z}}1
}
\renewcommand{\lstlistingname}{Listing}
\renewcommand{\lstlistlistingname}{Spis listingów}

% --- Literatura i spis treści ----------------------------------------
\usepackage[nottoc]{tocbibind}    % "Literatura" pojawi się w spisie treści
\addto\captionspolish{\renewcommand{\bibname}{Literatura}}
\usepackage{url}

\usepackage{tikz}
\usetikzlibrary{arrows.meta,positioning,calc,decorations.pathreplacing,shapes.geometric,fit,backgrounds,babel}
\usepackage{pgfplots}
\pgfplotsset{compat=1.18}
\definecolor{akcent}{RGB}{0,84,159}

\newenvironment{gdzie}{%
  \par\noindent gdzie:%
  \begin{itemize}[label={},leftmargin=2em,labelsep=0pt,itemsep=0pt,parsep=0pt,topsep=2pt]%
}{\end{itemize}\par}
\widowpenalty=10000
\clubpenalty=10000
\hyphenation{Mas-sa-chu-setts}
% --- Odnośniki klikalne w PDF (ładować jako ostatni) -----------------
\usepackage[hidelinks]{hyperref}
